\section{Introduction}

Agentic workflows (an LLM that plans, calls tools, delegates to sub-agents, and revises) are
evaluated end to end: run the whole workflow on a test set and score the outcome. That is slow and
expensive. A single public agent leaderboard spent about \$40{,}000 on 21{,}730 runs
\citep{@hal}, and cost is now reported next to accuracy because it decides which evaluations get
run at all \citep{@agentsthatmatter}. In practice, teams test few changes, on few tasks, and trust
noisy results.

Existing ways to cut this cost shrink something else. Benchmark compression keeps a few
informative test items \citep{@tinybenchmarks,@metabench,@efficientagents}. Simulators replace one
expensive component, such as the tools \citep{@toolemu,@stabletoolbench} or the user
\citep{@taubench}, often to build a new benchmark rather than to stand in for an existing system.
Prompt optimisers search on a cheaper model and carry the result to a stronger one
\citep{@cheapstrong,@gepa}. None of these asks the question a developer faces every day: \emph{can
I run a smaller version of my own workflow, and still make the same decisions about it?}

Engineers faced the same question a century ago with aircraft. A wind-tunnel model is not a small
aeroplane; it is valid because the few dimensionless quantities that govern the flow match those of
the real aircraft \citep{@buckingham}. This is the theory of \emph{similitude}. We ask what
similitude means for agent workflows: which parts of a workflow can be replayed, simulated, merged
or cut, and which properties must be kept, so that an evaluation of the miniature leads to the
same decisions as an evaluation of the original.

We make four contributions.
\begin{enumerate}
  \item \textbf{A problem and its yardstick.} We define \emph{evaluation-preserving reduction} and
  measure validity by decisions rather than scores: the same winner between variants, the same
  ranking, the same \emph{diagnosis} (the judge's reasons for failure) and \emph{fix transfer}
  (fixes found on the miniature improve the original). We read every miniature against a noise
  ceiling: the full workflow re-run with fresh samples.
  \item \textbf{A toolbox.} Thirteen reduction operators in three risk groups, from lossless
  prefix replay to fusing a whole pipeline into one LLM call with imagined tool results, and six
  rules for choosing among them.
  \item \textbf{A library} that applies the operators to unmodified workflow code through a single
  interception point, with a blind LLM judge that returns a score and a reason for every run.
  \item \textbf{A study} on three workflows (a single tool-calling agent on live Wikipedia, a
  customer-service agent on $\tau$-bench, and a five-role research pipeline) with \todo{N} variants
  each, including deliberately broken ones.
\end{enumerate}
\todo{one paragraph of headline findings once the results are in}
